\section{Matrix shells, fixed units and assembly}\label{sec:matrixshells}

\subsection{One energy unit for the whole N class}

Fix a class-N gap and the scale of~\eqref{eq:libscale}.
Use $\vartheta,\Omega,\Phi$ from~\eqref{eq:libinverse}--\eqref{eq:libfactor}.
The self-adjoint effective symbol is $h+\varepsilon M$, with
$\varepsilon=\delta/q$. If $m_0$ is the chosen scalar part of the
orbital matrix at the reference point, use the consistent convention
\begin{equation}\label{eq:libcenter}
 c'=c+\varepsilon m_0,\qquad
 W=z_\gamma(M-m_0I_2),\qquad y=Y+\varepsilon W.
\end{equation}
Then $y=z_\gamma(h+\varepsilon M-c')$.
For an energy $E$, put $Y_0=z_\gamma(E-c')$.
The operator $\operatorname{Op}(y)-Y_0$ is linear in $E$.
The endpoint-average to midpoint error is kept separately.

A single physical energy margin will be used for the entire gap class.
The reference size $N_{\rm ref}$ records the pair at the reference edge
and does not include the moving parameter $Y_0$. The factor $D$ pays
raw compression and bulk Schur errors. Depth enlarges the shell
threshold and penalty while $\epsilon_{\rm ref}$ and $A$ remain fixed.
Let $\overline H\ge1$ dominate the finite off-cell seminorms after
the fixed losses have been reserved. We may take
$\overline H\le e^{2\tau q}\mathcal R_q$.
Put
\begin{equation}\label{eq:libunit}
 \begin{split}
 N_{\rm ref}&=\sqrt{\mathrm{ex}}+1+\overline H,\\
 \epsilon_{\rm ref}&=(z_\gamma N_{\rm ref})^{-1},\\
 D&=C_D\{1+z_\gamma N_{\rm ref}
       (\gamma_{\rm cmp}+c_{\rm bulk}^2/m+\|W_{\rm bd}\|)\},\\
 A&=D\epsilon_{\rm ref},\qquad c_A=C_Wz_\gamma A=C_WD/N_{\rm ref}.
 \end{split}
\end{equation}
Here $C_D$ is a sufficiently large fixed constant for all assigned
bulk budgets, and $C_W$ is fixed below. The bulk coupling in this
definition is unweighted. The separate rank-two end coupling is
excluded from $D$. A boundary weight is retained when present.
This is one unit per gap and scale, independent of $E$.
It is finite at $\mathrm{ex}=0$.

The full class window is the contracted real chart with
\begin{equation}\label{eq:libfullwindow}
 |\zeta(E)|^2\le\mathrm{ex}+4\Lambda e^{-\xi q}.
\end{equation}
All assertions below hold for both branches and both signs of $Y_0$.
There is no restriction to shallow values of $|Y_0|$.

\subsection{The ordered product and the replacement shells}

On the full cell lattice let
\[
 \overline X=\operatorname{Op}(y)-Y_0,\qquad
 \overline V=\operatorname{Op}(y)+Y_0
             -2\operatorname{Op}(\vartheta(x)).
\]
Apply the partner after $\overline X$. The calculus of
Appendix~\ref{app:libcalculus} gives
\begin{equation}\label{eq:libproduct}
 \begin{split}
 \overline V\overline X&=\operatorname{Op}(p_1)+E_2,\\
 p_1&=\Phi(x)I_2+\varepsilon\{\widetilde Y,W\}_+
        +\varepsilon^2W^2
        +\frac{i\delta}{2}[y_k,y_t]
        +i\delta\{y,\vartheta\},\\
 \|E_2\|&\le\kappa\overline H\mathcal R_q=:e_2.
 \end{split}
\end{equation}
The braces with subscript $+$ denote the anticommutator; the last
braces are the ordered Poisson bracket. Both derivatives of
$\vartheta(x)$ are retained, including
$\partial_t\vartheta=\vartheta'(x)U'(t)$.
The first four terms of $p_1$ are Hermitian on real $k,t$.
The last term is anti-Hermitian. We use singular values and never
assert that the ordered product is self-adjoint.

Let $\mathcal D_j=(p_1)_0(t_j)$ be its exact zero Fourier block and
write $\mathcal D_j=\mathcal D_j^h+\mathcal D_j^a$.
Set
\[
 \mathcal D^h(t)=\overline\Phi(t)I_2+B(t),\qquad
 \overline\Phi(t)=\frac1{2\pi}\int_0^{2\pi}
             \Phi(\mathrm{ex}+T_s(k)+U(t))\,dk.
\]
The two ordered eigenvalues of $\mathcal D^h$ are denoted
$\lambda_-(t)\le\lambda_+(t)$.

\begin{lemma}[Block variation and global hopping]\label{lem:libhopping}
For the fixed derivative losses above,
\begin{align}
 \|B(t)\|&\le b_1\sqrt{\mathrm{ex}+4+U(t)}+b_2,
 &b_1&=C\sqrt\kappa\mathcal R_q,
 &b_2&=C(\sqrt\kappa+\kappa)\overline H\mathcal R_q,
 \label{eq:libB}\\
 \|B'(t)\|&\le C\sqrt\kappa\mathcal R_q
   (\sqrt{\mathrm{ex}+4+U(t)}+z_\gamma+\overline H)
       +C\kappa\mathcal R_q.\label{eq:libBder}
\end{align}
There is a core of $k_{\rm core}\le C\overline H\mathcal R_q$ cells
about the cell nearest $t_a$. On either side of the core each ordered
eigenvalue increases away from $t_a$, and
\begin{equation}\label{eq:libmatrixmono}
 \lambda_\ell(t)-\lambda_\ell(u)
     \ge\tfrac34(1-1/q)(U(t)-U(u))
\end{equation}
when $t$ is farther from $t_a$ than $u$ on that side.
For the full hopping kernel $\mathcal P_{ji}$ of
$\operatorname{Op}(p_1)$,
\begin{equation}\label{eq:libglobalhop}
 \sup_j\sum_{i\ne j}\|\mathcal P_{ji}\|,\quad
 \sup_i\sum_{j\ne i}\|\mathcal P_{ji}\|
 \le H_p:=C\{1+(\sqrt\kappa+\kappa)\overline H\mathcal R_q\}.
\end{equation}
The bound is independent of $\mathrm{ex}$ and $Y_0$.
\end{lemma}

\begin{proof}
The differentiated convolution estimates in
Appendix~\ref{app:libcalculus} bound $Y$ and $W$, with one extra
factor $q$ for $W$. Differentiating the dipole product uses
$\varepsilon(Y_tW+YW_t)$; differentiating the spin term uses
$y_{kt}$ and $y_{tt}$. The stated fixed seminorm orders control all
of them. They give~\eqref{eq:libB}--\eqref{eq:libBder}.
The anti-Hermitian block has a uniform bound $a_\infty$ with
\[
 a_\infty\le C\sqrt\kappa\overline H\mathcal R_q
                  +|\delta|\mathcal R_q/\rho_c.
\]
The last term is exponentially small for $b\ge2\xi$.

Weyl's min-max inequality, integrated along a side of the core,
gives
\[
 \lambda_\ell(t)-\lambda_\ell(u)
 \ge\overline\Phi(t)-\overline\Phi(u)-\int_u^t\|B'(v)\|\,|dv|.
\]
Since $|U'|\ge2\pi q\sqrt\Lambda\sqrt U$ on the chart,
each of the four ratios from~\eqref{eq:libBder} to
$(1-1/q)|U'|$ is at most $1/16$ outside an enlarged
$C\overline H\mathcal R_q$ core.
Here $\sqrt U$ changes by $\sqrt\kappa(1+o(1))$ per cell,
$\mathrm{ex}\le\mathrm{ex}_*\le\Lambda$, and
$z_\gamma\le q\sqrt\Lambda\mathcal R_q$.
Their sum is at most $1/4$, proving~\eqref{eq:libmatrixmono}.
No convexity or unique matrix-eigenvalue minimum is asserted inside
the core.

For~\eqref{eq:libglobalhop}, retain the zero Fourier matrix exactly.
Every nonzero Fourier coefficient of $\widetilde YW$ contains at
least one off-cell factor. The bounds
\[
 Y_{\rm off}\le\overline H\mathcal R_q,\quad
 W_{\rm off}\le q\overline H\mathcal R_q,\quad
 \|Y_0^{\rm Fourier}\|\le Cz_\gamma\mathcal R_q,\quad
 \|W_0^{\rm Fourier}\|\le qz_\gamma\mathcal R_q
\]
are bounds for absolute weighted Fourier sums. Thus its product
with $\varepsilon$ has row and column sums at most
$C|\delta|qz_\gamma\overline H\mathcal R_q
 \le C\sqrt\kappa\overline H\mathcal R_q$.
The $\varepsilon^2W^2$, spin and scalar-Poisson terms cost at most
$C\kappa\overline H\mathcal R_q$ and
$C\sqrt\kappa\overline H\mathcal R_q$.
Expand the scalar $\Phi$ in bounded $T_s(k)$; $\Phi'=1+O(1/q)$ and
the contracted-chart Cauchy bounds give off-cell sum $2+o(1)$.
Midpoint shifts and Fourier tails are included in these sums.
This proves the global estimate in both directions.
A large diagonal matrix has not been used as an off-cell bound.
\end{proof}

Keep the full diagonal matrix exactly. A cell belongs to the shell
only when one of its exact Hermitian block eigenvalues is close to
zero. The constant threshold pays hopping, non-Hermitian and
interface errors. Adding the diagonal bound again would overestimate
the width.

Fix a well $J$. Products of compressed factors retain an interface
norm $I_{\rm int}\le C\overline H(\overline H+h_\vartheta)$,
where $h_\vartheta$ is the off-cell sum of
$\operatorname{Op}(\vartheta)$. Choose a constant floor
\begin{equation}\label{eq:libWzero}
 \begin{split}
 W_0&\ge C_0(1+H_p+a_\infty+e_2+I_{\rm int}
                      +\kappa\mathcal R_q),\\
 W_0&\le C\mathcal R_q^C(1+\overline H^2+\kappa\overline H).
 \end{split}
\end{equation}
The second line follows by enlarging the fixed losses. Write
$n=|Y_0|+\sqrt{\mathrm{ex}}+1$ and define the replacement shell by
\begin{equation}\label{eq:libnewshell}
 \begin{split}
 W'&=\max\{W_0,\,2c_A(n+\overline H+1),\,4c_A^2\},\\
 \mathcal S&=\mathcal K_E\cup
 \{j\in J:\min_{\ell\in\{-,+\}}|\lambda_\ell(t_j)|<W'\}.
 \end{split}
\end{equation}
Here $\mathcal K_E$ is the whole core if the displayed inequality
holds somewhere in the core, and is empty otherwise.
In particular, the full diagonal majorant in~\eqref{eq:libB} is
absent from the threshold. It is already in the exact eigenvalues.

\begin{proposition}[Replacement shell geometry and penalized gap]
\label{prop:libshell}
For $2\xi\le b\le L-\xi$, large $q$, and every $E$ in
\eqref{eq:libfullwindow}, the shell has at most five components per
well, at most four if it misses the core. Each has width at most
\begin{equation}\label{eq:libshellwidth}
 C\left\{1+\sqrt{W_0/\kappa}
 +\frac{c_A(1+\overline H)}{\sqrt\kappa}
 +\sqrt{\frac{c_A(\sqrt{\mathrm{ex}}+1)}\kappa}\right\}
       +2k_{\rm core}.
\end{equation}
These widths obey~\eqref{eq:FPwidth}.
There is a bounded self-adjoint shell penalty, in energy units,
$M_E\un_{\mathcal S}$, such that the effective frame operator has
gap at least $4A$ at the window center and $3A$ for
$|E-E_c|\le A$ with the shell fixed at $E_c$. Its size satisfies
\begin{equation}\label{eq:libMbound}
 M=z_\gamma M_E\le C\{n+\sqrt{W_0}
                 +\overline H^2N_{\rm ref}/D+c_A\}.
\end{equation}
\end{proposition}

\begin{proof}
Put $\mathcal C=J\setminus\mathcal S$.
Off the core each ordered eigenvalue is monotone on each side and
the threshold is constant in the cell. Each branch gives at most
one interval on either side. Adding the conditional core gives the
component count.
For an interval with endpoint potential values $U_1\le U_2$,
\eqref{eq:libmatrixmono} gives $U_2-U_1\le CW'$.
Its cell count is at most
\[
 C\min\{\sqrt{W'},W'/\sqrt{U_2}\}/\sqrt\kappa+1.
\]
If $W'=W_0$ or $4c_A^2$, use the square-root bound.
For $W'=2c_A(n+\overline H+1)$ and $U_2\ge n^2/16$, use the second
bound. If $U_2<n^2/16$, the shell inequality and the full diagonal
bound~\eqref{eq:libB} imply
\[
 n\le C(\sqrt{\mathrm{ex}}+1+\sqrt{W'}+b_1+\sqrt{b_2}).
\]
The quantities $b_1^2,b_2$ are absorbed in $W_0$.
AM--GM applied to $c_A\sqrt{W'}$ proves
\eqref{eq:libshellwidth}.
No cell-dependent diagonal-size threshold has been added.

Here are both Schur sums for the off-shell gap. The floor gives
$d_j:=\sigma_{\min}(\mathcal D_j)\ge W'-a_\infty\ge0.99W'$ on
$\mathcal C$. For $\mathcal T_{ji}=\mathcal D_j^{-1}\mathcal P_{ji}$,
\begin{align*}
 \sup_j\sum_i\|\mathcal T_{ji}\|&\le H_p/(W'-a_\infty)\le1/32,\\
 \sup_i\sum_j\|\mathcal T_{ji}\|&\le
       (W'-a_\infty)^{-1}\sup_i\sum_j\|\mathcal P_{ji}\|\le1/32.
\end{align*}
Use operator norms, separately, for $E_2$ and the interface product:
$\|\mathcal D_{\mathcal C}^{-1}(E_{2,\mathcal C}
 -\mathrm{Int})\|\le1/32$.
Thus the total normalized remainder has norm at most $1/16$.
A norm-controlled remainder has not been assigned a row-sum bound.

Let $N$ be cellwise scalar, with
\[
 N_j=\sqrt{\mathrm{ex}+4+U_j}+|Y_0|+\overline H+\sqrt{W'}+1.
\]
The diagonal part of $N^{-1}\overline V_{\mathcal C\mathcal C}$
is bounded. Since $N_j\ge\overline H$, its off-cell row and column
sums satisfy, respectively,
\begin{align*}
 \sup_j\sum_{i\ne j}\frac{\|\overline V_{ji}\|}{N_j}
     &\le H_V/\overline H\le C,\\
 \sup_i\sum_{j\ne i}\frac{\|\overline V_{ji}\|}{N_j}
     &\le\overline H^{-1}\sup_i\sum_{j\ne i}
                       \|\overline V_{ji}\|\le C.
\end{align*}
Here the fixed seminorm losses have been absorbed into $\overline H$.
Hence its norm is at most a fixed $C_V$.
The product, with partner after $X$, now yields
\[
 C_V\|X_{\mathcal C\mathcal C}v\|
 \ge\|N^{-1}\mathcal D_{\mathcal C}(I+\mathcal T_{\rm total})v\|
 \ge\tfrac{15}{16}\min_j(d_j/N_j)\|v\|.
\]
This uses the order shown. It commutes neither $N$ nor
$\operatorname{Op}(\vartheta)$ past $X$.
For $U\le C(Y_0^2+\mathrm{ex}+W')$ the last minimum is at least
$cW'/(n+\sqrt{W'}+\overline H+1)\ge cc_A$.
In the outer region the scalar $\Phi$ dominates the matrix
majorant, the anti-Hermitian error and the floor, giving ratio at
least $c\sqrt U$, which is larger. Thus
$\|X_{\mathcal C\mathcal C}v\|\ge g\|v\|$ for the declared floor
$g=cc_A$.

Take $M=\|X_{\mathcal S\mathcal S}\|+g+3\overline H^2/g$.
Its shell block has gap $G\ge g+3\overline H^2/g$.
Diagonal scaling by $\operatorname{diag}(\sqrt G,\sqrt g)$ bounds
the normalized off-diagonal norm by $1/\sqrt3$. The whole penalized
operator therefore has singular gap at least
$(1-1/\sqrt3)g$. Increase $C_W$ once and divide by $z_\gamma$ to get
$4A$, with the endpoint error already paid in $D$.
Energy variation of size $A$ leaves $3A$.
If a block is empty, use the remaining gap or penalize the whole
finite well directly.
On shell cells the full diagonal bound implies
$\sqrt U\le C(n+\sqrt{W'}+\sqrt{W_0})$.
Substituting $g=cc_A$ gives~\eqref{eq:libMbound}.
In particular its $\overline H^2N_{\rm ref}/D$ term is retained.

Finally expand the unweighted definition of $D$ in
\eqref{eq:libunit}, using~\eqref{eq:libranktwoerrors} and
$c_{\rm bulk}^2/m\le\delta^2e^{\xi q/2}\mathcal R_q$.
Since $N_{\rm ref}\le\sqrt\Lambda\mathcal R_q$ and $r\le q/2$,
\begin{align*}
 \frac{c_A(1+\overline H)}{\sqrt\kappa}
       &\le\mathcal R_q^C(1+\kappa^{-1/2}),\\
 \frac{c_A(\sqrt{\mathrm{ex}}+1)}\kappa&\le CD/\kappa,\\
 \sqrt{D/\kappa}&\le e^{C\xi q}\mathcal R_q^C(1+\kappa^{-1/2}).
\end{align*}
For example the $\delta^2$ term in $D/\kappa$ is at most
$\mathcal R_q^C(N_{\rm ref}/\sqrt\Lambda)e^{\xi q/2}$.
The distance-dependent term is at most
$\mathcal R_q^CN_{\rm ref}e^{-(L-4\tau)(q-r)}/
 (\sqrt\Lambda|\delta|)$.
Its square root is bounded by the same type of slack times
$1+\kappa^{-1/2}$. The retained boundary weight and interwell tails
are smaller. Equations~\eqref{eq:libWzero} and
\eqref{eq:libshellwidth} now give~\eqref{eq:FPwidth}.
This calculation keeps the low-$b$ numerator $D$ throughout.
\end{proof}

\subsection{Weighted estimates at every depth}

\begin{proposition}[All-depth weighted budgets]\label{prop:libweighted}
Use the physical frame penalty synthesized from
Proposition~\ref{prop:libshell}. Let the weight vanish on shell cells,
jump at both-endpoint longer-arc bonds by at most
\[
 \mu=e^{-(\xi/2+10\tau)q},\qquad
 p_{\rm pad}=\lceil e^{(\xi/2+12\tau)q}\rceil,
\]
and cap it before the frame ends, keeping it constant in exterior
gaps. Throughout~\eqref{eq:libfullwindow}, the three weighted
budgets of Definition~\ref{def:FPstar} hold. All shells with their
padding lie inside the frame, and their fine-window count and
inverse widths are $e^{O(q)}$.
\end{proposition}

\begin{proof}
Let $F$ be~\eqref{eq:libF}. Physical weighted Schur estimates, with
both weight orientations, give for shell synthesis $U_{\mathcal S}$
\begin{equation}\label{eq:libonetails}
 \|(e^{\pm\varphi}-1)U_{\mathcal S}\|\le C\mu\sqrt F,
 \qquad
 \|V^*(e^{\pm\varphi}-1)U_{\mathcal S}\|\le C\mu F.
\end{equation}
The second bound uses two localized tails. Within a well the
weights grow more slowly than the frame tails. Across wells retain
$e^{\mu D_{\rm cap}}\mathcal R_qe^{-c/|\delta|}$, which is still
negligible. The proof is given in Appendix~\ref{app:liblocal}.

For open N,~\eqref{eq:libtunnel} and the Morse derivative comparison
imply
\begin{equation}\label{eq:libkeyproducts}
 \begin{split}
 z_\gamma\sqrt{\mathrm{ex}}F
       &\le e^{(d_h+6\tau)q}\mathcal R_q,\\
 \mathrm{ex}F&\le e^{(2d_h+4\tau)q-Lr}\mathcal R_q.
 \end{split}
\end{equation}
For closed or near-closed pairs use
$\sqrt{\mathrm{ex}}\le e^{(\xi/2+6\tau)q}$ and
$F\le\Lambda^{-1/2}e^{2\tau q}\mathcal R_q$ instead.
In either case
\begin{equation}\label{eq:libKref}
 K_{\rm ref}=z_\gamma N_{\rm ref}F
       \le e^{(\xi/2+8\tau)q}\mathcal R_q^C,\qquad
 \mu K_{\rm ref}=o(1).
\end{equation}
The physical weight perturbation has norm $O(\mu)$ and is supported
on the jump bonds. Its frame block is $O(\mu F)=o(A)$ by
\eqref{eq:libKref}. Its cross blocks are $O(\mu\sqrt F)$.
After squaring and dividing by $mA$, the extra factor is $C\mu/m$,
which tends to zero.

The penalty is nonlocal in physical space. Expanding its conjugation
with~\eqref{eq:libonetails} gives frame error $CM_E\mu F$, cross
errors $CM_E\mu\sqrt F$ and complement error $CM_E\mu^2F$.
The penalty is the only factor that grows with depth. Normalize its
frame error by $A=D/(z_\gamma N_{\rm ref})$; this leaves
$\mu MN_{\rm ref}F/D$. The deep part of $M$ is paid by
$z_\gamma\sqrt{\mathrm{ex}}F$, and its reference-mass part by
$\mathrm{ex}F$. Thus the same edge unit works at every depth.
More precisely, use the corrected penalty bound:
\begin{equation}\label{eq:libpenproduct}
 MN_{\rm ref}F\le e^{(8\tau+2d_h)q}\mathcal R_q^C.
\end{equation}
Here the depth-dependent term is paid by
\[
 \sqrt\Lambda e^{-\xi q/2}N_{\rm ref}F,
\]
which is bounded using the first line of~\eqref{eq:libkeyproducts}
or the near-closed cutoff. The squared reference term is paid by
$\mathrm{ex}F$ in its second line. From~\eqref{eq:libWzero},
$\sqrt{W_0}\le C\mathcal R_q^C
 (1+\overline H+\sqrt{\kappa\overline H})$.
This controls its term, and
$\overline H^2N_{\rm ref}^2F$ is controlled by the same
$\mathrm{ex}F$ bound. The last term is
$c_AN_{\rm ref}F=C_WDF$, with $D$ expanded in unweighted quantities
as in Proposition~\ref{prop:libshell}. This proves
\eqref{eq:libpenproduct}.
Consequently $\mu MN_{\rm ref}F=o(1)$.
The normalized penalty cross budget has an additional
$M_E\mu/m=o(1)$, since $M_E\le e^{4\tau q}\mathcal R_q^C$.
The complement physical error is $O(\mu)=o(m)$; its penalty error
has two one-tail factors and is $O(M_E\mu^2F)=o(m)$.
The adjoint cross block obeys the same bounds.

For location, a sufficient bound about the well bottom is
\begin{equation}\label{eq:libposition}
 |j-j_a|\le\mathcal R_q^C\left\{
 \sqrt{e/\kappa}+(\mathrm{ex}/\kappa)^{1/4}
 +\sqrt{W_0/\kappa}+c_A/\sqrt\kappa+k_{\rm core}\right\},
\end{equation}
where $e\le4\Lambda e^{-\xi q}$ bounds the spectral depth.
The extra orbital position term is at most $e^{(b/2+\text{chosen loss})q}$.
It is well inside the frame half-width
$e^{(b-\xi/4)q}/(2\pi q)$ for $b\ge2\xi$.
The other terms and $p_{\rm pad}$ are inside it as well.
This position bound does not count distance from the bottom as shell
width. Also $\mu p_{\rm pad}\ge e^{2\tau q}$.

Finally
\[
 e^{-(L+\text{chosen loss})q}\le\epsilon_{\rm ref}
      \le Cq^{-1/2}\Lambda^{-1/2}.
\]
The lower exponent is $L$, because $\mathrm{ex}$ can be large.
The class width is bounded, so there are $e^{O(q)}$ fine windows.
The weight and shell may be rebuilt in each smooth fine window.
This changes neither the fixed unit nor the absolute depth.
\end{proof}

\begin{remark}[Choice of energy unit]\label{rem:libvarying}
The available upper estimates do not prove the weighted budgets
with the varying unit $(z_\gamma\sqrt{\mathrm{ex}+e})^{-1}$ at all
depths. This is a limitation of that sufficient estimate. It is not
a theorem excluding every construction with varying units.
The fixed reference unit above is sufficient.
\end{remark}

\subsection{Assembly}

\begin{proof}[Proof of Theorem~\ref{thm:libration}]
We verify the five parts of Definition~\ref{def:FPstar}.
The finite units, covariant frames and physical kernels are given
by Lemma~\ref{lem:libframes}, Lemma~\ref{lem:libanchors} and
\eqref{eq:libunit}. Direct complement exclusion is
Proposition~\ref{prop:libcompression}. The unweighted definition
of $D$ pays the compression and bulk squared-coupling budgets.
It also gives $A=o(m)$, because
\[
 A\le C(\epsilon_{\rm ref}+\gamma_{\rm cmp}
                 +c_{\rm bulk}^2/m+\|W_{\rm bd}\|)=o(m).
\]
The shell gap and widths are Proposition~\ref{prop:libshell}; the
all-depth weighted budgets are Proposition~\ref{prop:libweighted}.
The localized three-block calculation in
Appendix~\ref{app:libboundary} supplies the retained end coupling.
W, top and bottom use the one-band package at the end of
Section~\ref{sec:innerframes}.

The full window~\eqref{eq:libfullwindow} covers both sides of an N
pair to $P$-depth $4\Lambda e^{-\xi q}$. The W and extreme windows
have the corresponding contracted depths. Their smooth cover
has $O(q)$ coarse classes and $e^{O(q)}$ fine windows.
It overlaps the high-$b$ plateau $\Lambda e^{-5\xi q/4}$ and the
target-dependent low-$b$ interior threshold in
Proposition~\ref{prop:FPtransport}. No window is reset to shallow
absolute depth. Thus all five parts hold, including closed pairs.

For any strict target $a>1-L/(2\beta)$, first fix the covering slack
as in Proposition~\ref{prop:FPtransport}. There are finitely many
constants in $e^{C\xi q}\mathcal R_q^C$. Choose $\xi$ inside that
slack, then $d_h$ after their largest derivative constant, and then
$q$ large. The factor $e^{2\tau q}$ and all other fixed exponential
losses are retained in this order. Proposition~\ref{prop:FPtransport}
now proves the stated tail estimate.
\end{proof}

\begin{remark}[Finite certificates; no computation in the proof]
\label{rem:libcertificates}
Independent interval calculations give the following enclosures
for sampled maxima of a frozen-fibre diagnostic. With
$\xi=L/4$, let $S_B$ be the maximum, over the two real Bloch corners,
of the sum of pair-projector diagonals at both jump-bond endpoints,
and put $Q_{\rm old}=z_\gamma\sqrt{\mathrm{ex}+4\Lambda e^{-\xi q}}S_B$.
For the finite proxy class $\mathrm{ex}\le\Lambda/q^3$, both turning
signs, the certified intervals for the maximum are
\[
\begin{array}{c|c|c|c}
 \lambda&p&q&\max Q_{\rm old}\\ \hline
 2&8&13&[65.567943,65.567944]\\
 2&13&21&[917.687046,917.687047]\\
 2&21&34&[56263.586146,56263.586147]\\
 3&21&34&[6992639.730639,6992639.730640]\\
 7&8&17&[280360.523174,280360.523175]
\end{array}
\]
These are finite certificates, rounded outward from rational
interval endpoints. The central touching pair at even $q$ is
certified as a rank-two projector using the exact touching identity,
residual and Gram bounds. The proxy class is different from
\eqref{eq:libclasses}; two Bloch corners do not certify a uniform
physical operator estimate. The intervals do not prove an
asymptotic obstruction for varying units. None is used above.
\end{remark}
